\documentclass[preprint,11pt,authoryear]{elsarticle}

%─── Packages ────────────────────────────────────────────────────────────────
\usepackage[T1]{fontenc}
\usepackage{lmodern}
\usepackage{amsmath,amssymb,amsthm,bm,bbm}
\usepackage[dvipsnames]{xcolor}
\usepackage{booktabs,colortbl,array,multirow,threeparttable}
\DeclareMathOperator*{\argmin}{arg\,min}
\usepackage{graphicx,float,caption,subcaption}
\graphicspath{{figures/}}
\usepackage{algorithm,algpseudocode}
\usepackage{natbib}
\usepackage[hidelinks]{hyperref}
\hypersetup{colorlinks=false, allbordercolors={1 1 1}, pdfborder={0 0 0},
  pdftitle={Online Supplement to: Conformal Recalibration of Tail Quantile Forecasts under Temporal Dependence}}
\usepackage{placeins}
\usepackage{setspace,microtype,enumitem}
\onehalfspacing

%─── Colours ─────────────────────────────────────────────────────────────────
\definecolor{pastelgreen}{rgb}{0.0,0.6,0.0}

%─── Revision markup ─────────────────────────────────────────────────────────
% Retained as a no-op so that generated fragments emitting \revised still build.
\newcommand{\revised}[1]{#1}

%─── Generated figures ───────────────────────────────────────────────────────
% The same macro file the manuscript uses, so a number quoted in the supplement
% cannot drift from the same number in the main text.
% Generated by scripts/paper_numbers.py -- do not edit by hand.
% Every figure the manuscript asserts, recomputed from the artefacts.
% Macro names carry their panel: Main = 16 x 24 = 384 pairs (Table 1);
% Seq = 13 x 24 = 312 pairs, the forecasters with stored series.

\newcommand{\nMainForecasters}{16}
\newcommand{\nMainAssets}{24}
\newcommand{\nMainPairs}{384}
\newcommand{\nSeqForecasters}{13}
\newcommand{\nSeqPairs}{312}
\newcommand{\nMainGreen}{335}
\newcommand{\nMainGreenPct}{87.2}
\newcommand{\nMainKupiecRawPasses}{123}
\newcommand{\nMainKupiecCorPasses}{273}
\newcommand{\nMainBestKupiec}{15}
\newcommand{\nMainBestKupiecModel}{CAViaR-AS}
\newcommand{\nMainCCDefined}{180}
\newcommand{\nMainCCUndefined}{204}
\newcommand{\nMainCCUndefPct}{53.1}
\newcommand{\nMainCCPass}{97}
\newcommand{\nMainCCPassPct}{53.9}
\newcommand{\nMainCCAsPassPct}{78.4}
\newcommand{\nMainRawPiMin}{0.0110}
\newcommand{\nMainRawPiMax}{0.0333}
\newcommand{\nMainRMin}{0.001}
\newcommand{\nMainRMax}{0.357}
\newcommand{\nMainRTruncOne}{17.3}
\newcommand{\nMainRTruncTwo}{23.5}
\newcommand{\nRawPiTimesFM}{0.0143}
\newcommand{\nRawPiMoiraiOne}{0.0154}
\newcommand{\nRawPiMoiraiTwo}{0.0178}
\newcommand{\nRawPiLagLlama}{0.0294}
\newcommand{\nRawPiChronosSmallA}{0.0173}
\newcommand{\nRawPiChronosMiniA}{0.0177}
\newcommand{\nRawPiGJR}{0.0200}
\newcommand{\nRawPiGJRt}{0.0145}
\newcommand{\nRawPiCAViaRAS}{0.0110}
\newcommand{\nWidthRatioTimesFM}{1.016}
\newcommand{\nWidthRatioMoiraiTwo}{1.001}
\newcommand{\nGJRnegqV}{0}
\newcommand{\nSeqKupiecRejRawOne}{69.9}
\newcommand{\nSeqCCUndefRawOne}{34.6}
\newcommand{\nSeqCCRejRawOne}{27.5}
\newcommand{\nSeqCCUndefCorOne}{53.5}
\newcommand{\nSeqCCRejCorOne}{48.3}
\newcommand{\nSeqKupiecRejRawTwoFive}{62.8}
\newcommand{\nSeqCCUndefRawTwoFive}{11.9}
\newcommand{\nSeqCCRejRawTwoFive}{19.3}
\newcommand{\nSeqCCUndefCorTwoFive}{26.6}
\newcommand{\nSeqCCRejCorTwoFive}{30.6}
\newcommand{\nSeqKupiecRejRawFive}{47.8}
\newcommand{\nSeqCCUndefRawFive}{1.0}
\newcommand{\nSeqCCRejRawFive}{22.0}
\newcommand{\nSeqCCUndefCorFive}{4.5}
\newcommand{\nSeqCCRejCorFive}{24.5}
\newcommand{\nSeqKupiecRejRawTen}{40.1}
\newcommand{\nSeqCCUndefRawTen}{0.0}
\newcommand{\nSeqCCRejRawTen}{26.0}
\newcommand{\nSeqCCUndefCorTen}{0.0}
\newcommand{\nSeqCCRejCorTen}{28.5}
\newcommand{\nTruncCorPi}{0.0108}
\newcommand{\nTruncCorGreen}{19}
\newcommand{\nGateStaticApplied}{259}
\newcommand{\nGateStaticSkipped}{53}
\newcommand{\nGateStaticAvoided}{17}
\newcommand{\nGateStaticForgone}{36}
\newcommand{\nGateStaticUpgrades}{181}
\newcommand{\nGateStaticUpgradesKept}{171}
\newcommand{\nGateStaticOracleApplied}{218}
\newcommand{\nGateStaticOracleSkipped}{94}
\newcommand{\nGateStaticOracleAvoided}{58}
\newcommand{\nGateStaticOracleForgone}{36}
\newcommand{\nGateStaticOracleUpgrades}{181}
\newcommand{\nGateStaticOracleUpgradesKept}{181}
\newcommand{\nDegradedStatic}{66}
\newcommand{\nDegradedStaticZoneUp}{6}
\newcommand{\nDegradedStaticNoChange}{60}
\newcommand{\nDegradedStaticAlreadyGreen}{59}
\newcommand{\nDegradedStaticSameNotGreen}{1}
\newcommand{\nDegradedStaticZoneDown}{0}
\newcommand{\nWellCalStaticN}{94}
\newcommand{\nWellCalStaticWorse}{58}
\newcommand{\nWellCalStaticMeanPct}{-1.3}
\newcommand{\nWellCalStaticWilcoxon}{8.9 \times 10^{-5}}
\newcommand{\nGateRollApplied}{259}
\newcommand{\nGateRollSkipped}{53}
\newcommand{\nGateRollAvoided}{44}
\newcommand{\nGateRollForgone}{9}
\newcommand{\nGateRollUpgrades}{204}
\newcommand{\nGateRollUpgradesKept}{184}
\newcommand{\nGateRollOracleApplied}{218}
\newcommand{\nGateRollOracleSkipped}{94}
\newcommand{\nGateRollOracleAvoided}{89}
\newcommand{\nGateRollOracleForgone}{5}
\newcommand{\nGateRollOracleUpgrades}{204}
\newcommand{\nGateRollOracleUpgradesKept}{204}
\newcommand{\nDegradedRoll}{174}
\newcommand{\nDegradedRollZoneUp}{71}
\newcommand{\nDegradedRollNoChange}{103}
\newcommand{\nDegradedRollAlreadyGreen}{100}
\newcommand{\nDegradedRollSameNotGreen}{1}
\newcommand{\nDegradedRollZoneDown}{2}
\newcommand{\nWellCalRollN}{94}
\newcommand{\nWellCalRollWorse}{89}
\newcommand{\nWellCalRollMeanPct}{-7.4}
\newcommand{\nWellCalRollWilcoxon}{10^{-16}}
\newcommand{\nGateRollAvoidedPct}{25.3}
\newcommand{\nGateRollOracleAvoidedPct}{51.1}
\newcommand{\nGateStaticAvoidedPct}{25.8}
\newcommand{\nGateRollUpgradesLost}{20}
\newcommand{\nGateStaticUpgradesLost}{10}
\newcommand{\nDegradedRollNoChangePct}{59.2}
\newcommand{\nDispDefaultSmall}{0.121}
\newcommand{\nDispDefaultMini}{0.115}
\newcommand{\nDispTwoHundredSmall}{0.331}
\newcommand{\nDispTwoHundredMini}{0.328}
\newcommand{\nDispThousandSmall}{0.958}
\newcommand{\nDispThousandMini}{0.945}
\newcommand{\nDispFullSmall}{1.087}
\newcommand{\nDispFullMini}{1.059}
\newcommand{\nDispTempLowSmall}{0.120}
\newcommand{\nDispTempLowMini}{0.114}
\newcommand{\nDispTempHighSmall}{0.121}
\newcommand{\nDispTempHighMini}{0.115}
\newcommand{\nDispNucleusLowSmall}{0.115}
\newcommand{\nDispNucleusLowMini}{0.109}
\newcommand{\nDoseCells}{1600}
\newcommand{\nPiSmallDefaultOne}{0.3884}
\newcommand{\nKupSmallDefaultOne}{0}
\newcommand{\nPiSmallDefaultTen}{0.4177}
\newcommand{\nKupSmallDefaultTen}{0}
\newcommand{\nRatioSmallDefaultOne}{38.837}
\newcommand{\nRatioSmallDefaultTen}{4.177}
\newcommand{\nPiSmallAnalyticOne}{0.0173}
\newcommand{\nKupSmallAnalyticOne}{8}
\newcommand{\nPiSmallAnalyticTen}{0.1027}
\newcommand{\nKupSmallAnalyticTen}{20}
\newcommand{\nRatioSmallAnalyticOne}{1.733}
\newcommand{\nRatioSmallAnalyticTen}{1.027}
\newcommand{\nPiMiniDefaultOne}{0.4188}
\newcommand{\nKupMiniDefaultOne}{0}
\newcommand{\nPiMiniDefaultTen}{0.4477}
\newcommand{\nKupMiniDefaultTen}{0}
\newcommand{\nRatioMiniDefaultOne}{41.884}
\newcommand{\nRatioMiniDefaultTen}{4.477}
\newcommand{\nPiMiniAnalyticOne}{0.0177}
\newcommand{\nKupMiniAnalyticOne}{8}
\newcommand{\nPiMiniAnalyticTen}{0.0981}
\newcommand{\nKupMiniAnalyticTen}{20}
\newcommand{\nRatioMiniAnalyticOne}{1.772}
\newcommand{\nRatioMiniAnalyticTen}{0.981}
\newcommand{\nGateBlocked}{4}
\newcommand{\nGateSeries}{13}
\newcommand{\nGateExtremesSmall}{19}
\newcommand{\nGateExtremesMini}{15}
\newcommand{\nWCDMPairs}{30}
\newcommand{\nWCDMSigBoot}{30}
\newcommand{\nWCDMSigAsymp}{29}
\newcommand{\nWCKupAsympGJR}{0.029}
\newcommand{\nWCKupBootGJR}{0.234}
\newcommand{\nLagForecasters}{13}
\newcommand{\nLagImmediate}{7}
\newcommand{\nLagMidDays}{77}
\newcommand{\nLagMidN}{4}
\newcommand{\nLagLateDays}{161}
\newcommand{\nLagLateN}{2}
\newcommand{\nLitClosureRMin}{0.005}
\newcommand{\nLitClosureRMax}{1.70}
\newcommand{\nGapDeltaMoment}{0.046}
\newcommand{\nGapUndGarchT}{10.6}
\newcommand{\nLitAnalyticDates}{40}
\newcommand{\nLitAnalyticSdPct}{0.55}
\newcommand{\nLitAnalyticGrid}{0.025}
\newcommand{\nRawPiGJRDefective}{0.0041}
\newcommand{\nRawPiGJRDefectivePct}{0.4}
\newcommand{\nPoolPiEWMA}{0.0119}
\newcommand{\nPoolPiMoiraiTwo}{0.0114}
\newcommand{\nPoolPiSmallAnalytic}{0.0113}
\newcommand{\nPoolPiGJR}{0.0111}
\newcommand{\nPoolBootEWMA}{0.058}
\newcommand{\nPoolClusterEWMA}{0.035}
\newcommand{\nuMCBShareWellSpec}{0.270}
\newcommand{\nuMCBShareAll}{0.329}
\newcommand{\nuMCBSharePct}{27.0}
\newcommand{\nLitTempGap}{0.001}
\newcommand{\nLitNucleusGap}{0.005}
\newcommand{\nLitSupportPct}{98.8}
\newcommand{\nLitRatioSmallTwoFive}{1.34}
\newcommand{\nLitRatioSmallFive}{1.14}
\newcommand{\nLitRatioMiniTwoFive}{1.32}
\newcommand{\nLitRatioMiniFive}{1.11}
\newcommand{\nLitAlphaResp}{1.08}
\newcommand{\nLitMainRMinAll}{0.001}
\newcommand{\nLitMainRMaxOK}{0.357}
\newcommand{\nLitMainROkN}{14}
\newcommand{\nLitQSMin}{4.67}
\newcommand{\nLitQSMax}{5.83}
\newcommand{\nLitWMin}{0.98}
\newcommand{\nLitWMax}{1.14}
\newcommand{\nLitCorPiTypical}{0.011}
\newcommand{\nLitMoiraiGapPP}{0.24}
\newcommand{\nLitLambdaTail}{1.9\times 10^{-7}}
\newcommand{\nLitQSDefaultSmall}{5.81}
\newcommand{\nLitQSAnalyticSmall}{5.03}
\newcommand{\nLitQSDefaultMini}{5.69}
\newcommand{\nLitQSAnalyticMini}{5.00}
\newcommand{\nSeqDQCells}{312}
\newcommand{\nSeqDQRejRaw}{81.7}
\newcommand{\nSeqDQRejCor}{53.8}
\newcommand{\nGapDeltaFree}{0.218}
\newcommand{\nGapUndFree}{75.5}
\newcommand{\nGapDeltaUni}{0.066}
\newcommand{\nGapUndUni}{49.4}
\newcommand{\nGapDeltaGateNow}{0.024}
\newcommand{\nGapUndGateNow}{30.9}
\newcommand{\nGapBandNow}{-1.800}
\newcommand{\nGapBandLower}{-3.500}
\newcommand{\nGapCellGoodWorst}{-1.947}
\newcommand{\nGapCellGoodBest}{-2.735}
\newcommand{\nGapCellTruncBest}{-0.283}
\newcommand{\nGapCellTruncWorst}{-0.074}
\newcommand{\nGapGapWidth}{1.664}
\newcommand{\nGapMarginUpper}{0.147}
\newcommand{\nGapMarginLower}{0.765}
\newcommand{\nGapCellsLowerBlocks}{0}
\newcommand{\nGapCellsTotal}{312}
\newcommand{\nGapCellsTrunc}{48}
\newcommand{\nGapDeltaChronos}{0.388}
\newcommand{\nGapQTrueStd}{-2.606}
\newcommand{\nGapUndTight}{25.6}
\newcommand{\nGapMarginTight}{0.007}
\newcommand{\nGapBandTight}{-1.940}
\newcommand{\nGapDMt}{0.399}
\newcommand{\nGapDMp}{0.69}
\newcommand{\nGapQSGapPct}{0.008}
\newcommand{\nGapVaRHonest}{2.61}
\newcommand{\nGapVaRAlt}{1.32}
\newcommand{\nGapRhoLo}{0.18}
\newcommand{\nGapRhoHi}{0.67}
\newcommand{\nGapDeltaHatLo}{0.109}
\newcommand{\nGapDeltaHatHi}{0.138}
\newcommand{\nGapEmpCoverage}{98.9}
\newcommand{\nGapGapAblCovid}{0.0058}
\newcommand{\nGapGapAblFull}{0.0005}
\newcommand{\nGapFisherKupiec}{0.49}
\newcommand{\nGapFisherSevere}{0.00078}
\newcommand{\nGapAblRhoLo}{+0.05}
\newcommand{\nGapAblRhoHi}{+0.46}
\newcommand{\nGapAblPairs}{4}
\newcommand{\nTLTau}{0.016}
\newcommand{\nTLTauOverAlpha}{1.6}
\newcommand{\nTLUndTFive}{12.4}
\newcommand{\nTLUndTThree}{16.3}
\newcommand{\nTLUndNormal}{7.8}
\newcommand{\nMCGreenTFiveSmall}{61.6}
\newcommand{\nMCGreenTFiveLarge}{74.2}
\newcommand{\nMCPiTFive}{0.0148}
\newcommand{\nMCGreenSkewSmall}{36.4}
\newcommand{\nMCGreenSkewLarge}{0.4}
\newcommand{\nMCTMax}{10{,}000}
\newcommand{\nMCTMin}{500}
\newcommand{\nMCReps}{500}
\newcommand{\nMCDgps}{5}
\newcommand{\nGateRollUpgradesLostWorse}{11}
\newcommand{\nGateRollUpgradesLostNet}{9}
\newcommand{\nGateRollUpgradeAndDeteriorate}{71}
\newcommand{\nSpearmanRPi}{0.991}
\newcommand{\nSpearmanRPiN}{14}


%─── Theorem environments ────────────────────────────────────────────────────
\newtheorem{theorem}{Theorem}[section]
\newtheorem{corollary}[theorem]{Corollary}
\newtheorem{definition}[theorem]{Definition}
\newtheorem{remark}[theorem]{Remark}
\newtheorem{assumption}[theorem]{Assumption}
\newtheorem{proposition}[theorem]{Proposition}
\newtheorem{lemma}[theorem]{Lemma}

%─── Notation ────────────────────────────────────────────────────────────────
\newcommand{\VaR}{\mathrm{VaR}}
\newcommand{\ES}{\mathrm{ES}}
\newcommand{\E}{\mathbb{E}}
\renewcommand{\P}{\mathsf{P}}
\newcommand{\1}{\boldsymbol{1}}
\newcommand{\qVstat}{\hat{q}_{\scriptscriptstyle\mathrm{V}}^{\,\mathrm{stat}}}
% Backward-compatible alias: three robustness fragments predate the
% \qV -> \qVstat rename and still emit \qV. Aliasing here avoids editing
% generated files, which would be overwritten on the next pipeline run.
\providecommand{\qV}{\qVstat}
\newcommand{\qVroll}{\hat{q}_{{\scriptscriptstyle\mathrm{V}},t}^{\,\mathrm{roll}}}
\newcommand{\qE}{\hat{q}_{\scriptscriptstyle\mathrm{E}}}
\newcommand{\EScqr}{\widehat{\ES}^{\,\mathrm{cqr}}}
\newcommand{\CI}{\hat{C}}
\newcommand{\Fhat}{\hat{\mathcal{F}}}
\newcommand{\quantlet}[1]{\href{#1}{\raisebox{-0.15em}{\includegraphics[height=0.9em]{ql_logo.png}}}}
\newcommand{\quantinar}[1]{\href{#1}{\raisebox{-0.15em}{\includegraphics[height=0.9em]{qr_logo.png}}}}

%══════════════════════════════════════════════════════════════════════════════

%══════════════════════════════════════════════════════════════════════════════
%  ONLINE SUPPLEMENT to
%  "What Backtests Cannot Diagnose: Structural Validation of Tail-Risk
%  Forecasting Pipelines"
%
%  Sections moved out of the manuscript's appendices on 2026-08-20 to keep the
%  submitted paper to a readable length. Numbering is S.1-S.4; the manuscript
%  refers to them by those numbers rather than by \ref, so neither document
%  needs the other to compile.
%══════════════════════════════════════════════════════════════════════════════


\begin{document}

\title{Proofs for \emph{Recalibrating Tail-Event Forecasts under Temporal Dependence}}
\date{}
\maketitle

\noindent This document carries the proofs of the results stated in the
manuscript. It is part of the replication material rather than of the
submitted article, and it is referenced from the text by name.
Theorem~4.5 is a one-sided specialization of a published result and its
proof is included for completeness; the attribution is in the manuscript.

\section{Proofs}
\label{app:proofs}
%============================================================

Section~4 of the manuscript states the assumption, the two lemmas, the coverage
theorem, the GARCH corollary and the two rolling-estimator bounds, together with
their scope. This appendix carries the proofs, and each is labelled by the
numbering of the statement it proves. No statement is repeated here in a second
form.

For reference while reading them, Assumption~4.1 of the manuscript requires
(A1) that the score process $S_t = \hat q_t^{lo} - r_t$ be strictly stationary
and geometrically $\beta$-mixing, $\beta(k)\le C\rho^k$; (A2) that the marginal
law $F_S$ have a continuous density bounded away from zero near
$q_{1-\alpha}$; and (A3) that the evaluated score be separated from the
calibration block by a gap $g_n\to\infty$ with $g_n/n\to0$.

\subsection{Proof of Lemma~4.2 (inheritance of weak dependence)}
\label{app:lem_score_mixing}

\begin{remark}[Why the context must be finite]
	\label{rem:finite_context}
	The finiteness of $m$ is not a convenience. Allowing
	$\hat q_t^{lo}=g(r_{t-1},r_{t-2},\ldots)$ with an infinite context does not
	yield $\beta_S(k)\le\beta_r(k-1)$: the proof requires
	$\sigma(S_s : s\ge t+k)\subseteq\sigma(r_u : u\ge t+k-1)$, and with an
	infinite context $S_s$ depends on returns arbitrarily far before $t+k-1$,
	so the inclusion fails and no separation of the two $\sigma$-fields is
	obtained. Restricting $g$ to a window of length $m$ is what makes the
	blocking argument available, and for all but one forecaster here the
	restriction holds by construction: $m = 512$ for the foundation models and
	$m = 250$ for the GARCH-family and historical-simulation benchmarks
	(Section~3.3.3 of the manuscript).

	EWMA is the exception, and it is stated rather than rounded off. Its
	conditional variance is the RiskMetrics recursion, whose memory is infinite,
	so the lemma does not apply to it as literally implemented; the manuscript
	reconciles this wherever the 250-observation convention is quoted. What can be
	said is quantitative: at $\lambda = 0.94$ the weight beyond lag $k$ is
	$\lambda^k$, which is $1.9\times10^{-7}$ at $k = 250$ and falls below
	double-precision machine epsilon at $k = 583$. The lemma applies to the filter
	truncated at $m = 600$, which is indistinguishable from the recursion in the
	arithmetic that produced the series. That is an approximation argument rather
	than an exemption, and it is the only one in the paper.
\end{remark}

\begin{proof}
	By assumption, $\{r_t\}$ is strictly stationary. Since $\hat q_t^{lo}$ is
	a time-invariant measurable functional of the past, the joint process
	$(\hat q_t^{lo},r_t)$ is also strictly stationary. Therefore,
	$S_t=\hat q_t^{lo}-r_t$ inherits strict stationarity.
	
	Let $\mathcal S_{-\infty}^t=\sigma(S_s:s\le t)$ and
	$\mathcal S_{t+k}^{\infty}=\sigma(S_s:s\ge t+k)$. By
	equation~(10) of the manuscript, $S_t$ is a measurable function of
	$(r_t,r_{t-1},\ldots,r_{t-m})$, so
	\[
	\mathcal S_{-\infty}^t
	\subseteq
	\mathcal R_{-\infty}^t:=\sigma(r_s:s\le t),
	\]
	and, crucially, the forward field reaches back only $m$ steps:
	\[
	\mathcal S_{t+k}^{\infty}
	\subseteq
	\mathcal R_{t+k-m}^{\infty}:=\sigma(r_s:s\ge t+k-m).
	\]
	It is this second inclusion that fails without a finite context
	(Remark~S.9.1).
	By monotonicity of $\beta$-mixing coefficients under measurable
	transformations,
	\[
	\beta_S(k)
	\le
	\beta(\mathcal R_{-\infty}^t,\mathcal R_{t+k-m}^{\infty})
	=
	\beta_r(k-m),
	\qquad k>m,
	\]
	which is equation~(11) of the manuscript.
\end{proof}

\subsection{Proof of Lemma~4.4 (empirical quantile deviation)}
\label{app:lem_quantile}

\begin{proof}
	Under geometric $\beta$-mixing as in Assumption~4.1(A1) of the manuscript,
	the empirical CDF satisfies the uniform concentration bound
	\begin{equation}
		\label{eq:app_ecdf_bound}
		\sup_{s\in\mathbb R}|\hat F_n(s)-F_S(s)|
		=
		\mathcal O_p\!\left(\sqrt{\frac{\log n}{n}}\right);
	\end{equation}
	see \citet[Theorem 2.1]{yu1994rates} and
	\citet[Theorem 7.2]{rio2017asymptotic}.
	 Since $\qVstat$ is the
	empirical $(1-\alpha)$-quantile,
	\[
	\hat F_n(\qVstat)=1-\alpha+\mathcal O(1/n).
	\]
	Thus,
	\[
	|F_S(\qVstat)-(1-\alpha)|
	\le
	\sup_s|\hat F_n(s)-F_S(s)|+\mathcal O(1/n).
	\]
	By the density condition in Assumption~4.1 of the manuscript,
	$f_S(q_{1-\alpha})$ is bounded away from zero in a neighbourhood of
	$q_{1-\alpha}$. A first-order expansion around $q_{1-\alpha}$ gives
	\[
	|\qVstat-q_{1-\alpha}|
	\le
	\frac{1}{f_S(q_{1-\alpha})}
	\left[
	\sup_s|\hat F_n(s)-F_S(s)|+\mathcal O(1/n)
	\right],
	\]
	which yields the stated rate.
\end{proof}

\subsection{Proof of Theorem~4.5 (coverage under $\beta$-mixing)}
\label{app:thm_coverage}

\begin{proof}
	The proof proceeds in three steps.
	
	\textit{Step 1: Reduction to score control.}
	By definition of the nonconformity score
	$S_t=\hat q_t^{lo}-r_t$,
	\begin{equation}
		\label{eq:proof_score_event}
		S_{n+g_n+1}\le\qVstat
		\Longleftrightarrow
		r_{n+g_n+1}\ge \hat q_{n+g_n+1}^{lo}-\qVstat.
	\end{equation}
	It therefore suffices to show
	\begin{equation}
		\label{eq:proof_target_coverage}
		\mathbb P(S_{n+g_n+1}\le\qVstat)
		\ge
		1-\alpha-\Delta_n.
	\end{equation}
	
\textit{Step 2: Quantile approximation.}
Let $q_{1-\alpha}=F_S^{-1}(1-\alpha)$. By
Lemma~4.4 of the manuscript,
\begin{equation}
	\label{eq:proof_qv_rate}
	|\qVstat-q_{1-\alpha}|
	=
	\mathcal O_p\!\left(\sqrt{\frac{\log n}{n}}\right).
\end{equation}
By the density condition,
\begin{equation}
	\label{eq:proof_fs_qv_bound}
	F_S(\qVstat)
	\ge
	1-\alpha
	-
	C\sqrt{\frac{\log n}{n}}
\end{equation}
for a finite constant $C>0$.

\textit{Step 3: Gap-separated dependence control.}
By $\beta$-mixing, conditional on the calibration sample
$\mathcal D_n=(S_1,\ldots,S_n)$,
\begin{equation}
	\label{eq:proof_conditional_gap}
	\bigl|
	\mathbb P(S_{n+g_n+1}\le\qVstat\mid\mathcal D_n)
	-
	F_S(\qVstat)
	\bigr|
	\le
	\beta(g_n).
\end{equation}
Taking expectations over $\mathcal D_n$ and combining with
\eqref{eq:proof_fs_qv_bound} gives
\[
\mathbb P(S_{n+g_n+1}\le\qVstat)
\ge
\mathbb E\!\left[F_S(\qVstat)\right]-\beta(g_n)
\ge
1-\alpha
-
C\sqrt{\frac{\log n}{n}}
-\beta(g_n).
\]
Setting
\begin{equation}
	\label{eq:proof_delta_n}
	\Delta_n
	=
	C\sqrt{\frac{\log n}{n}}+\beta(g_n)
\end{equation}
completes the proof.
\end{proof}

\subsection{Proof of Corollary~4.6 (GARCH data-generating processes)}
\label{app:cor_garch}

\begin{proof}
	Under $\omega>0$, $\alpha_1,\beta_1\ge0$, and
	$\mathbb E[\log(\alpha_1\varepsilon_t^2+\beta_1)]<0$, the GARCH(1,1)
	process admits a unique strictly stationary solution
	\citep[Theorem~2.10]{francq2019garch}. Under standard regularity
	conditions, it is geometrically $\beta$-mixing
	\citep{francq2019garch,lindner2009stationarity}:
	\[
	\beta_r(k)\le a\rho^k,
	\qquad \rho\in(0,1).
	\]
	By Lemma~4.2 of the manuscript, the score process is also strictly
	stationary and geometrically $\beta$-mixing. The density condition holds
	under the assumed continuous innovation density near the target quantile.
	
	Choosing $g_n=c\log n$ with $c>1/|\log\rho|$ gives
	\begin{equation}
		\label{eq:app_garch_gap_rate}
		\beta(g_n)\le a\rho^{c\log n}
		=
		\mathcal O(n^{-c|\log\rho|})
		=
		o\!\left(\sqrt{\frac{\log n}{n}}\right).
	\end{equation}
	Theorem~4.5 of the manuscript therefore applies with
	\[
	\Delta_n
	=
	\mathcal O\!\left(\sqrt{\frac{\log n}{n}}\right).
	\]
\end{proof}

\subsection{Proof of Proposition~4.7 (online quantile tracking)}
\label{app:ogd_regret}

Proposition~4.7 of the manuscript is deterministic: it holds for every score
sequence, with no probabilistic assumption whatsoever. Its scope relative to the
estimator actually used in the empirical sections is stated after the proof.


\begin{proof}[Proof of Proposition~4.7 of the manuscript]
	\emph{Step 1: the iterates stay bounded, without projection.} We claim
	$q_t\in[-R-\eta\alpha,\;R+\eta(1-\alpha)]$ for every $t\ge1$. The claim
	holds at $t=1$ by assumption. Suppose it holds at $t$. For the upper end,
	either $q_t\le R$, in which case
	$q_{t+1}\le q_t+\eta(1-\alpha)\le R+\eta(1-\alpha)$; or $q_t>R\ge s_t$, in
	which case $\mathbf 1\{s_t>q_t\}=0$, the update is $-\eta\alpha<0$, and
	$q_{t+1}<q_t$. For the lower end, either $q_t\ge -R$, so
	$q_{t+1}\ge q_t-\eta\alpha\ge -R-\eta\alpha$; or $q_t<-R\le s_t$, so
	$\mathbf 1\{s_t>q_t\}=1$, the update is $+\eta(1-\alpha)>0$, and
	$q_{t+1}>q_t$. In both directions the iterate is pulled back as soon as it
	leaves $[-R,R]$, so no projection step is needed and the range has width at
	most $2R+\eta$.

	\emph{Step 2: telescoping.} Summing the second form of
	equation~(15) of the manuscript over $t=1,\ldots,T$,
	\[
	\eta\sum_{t=1}^{T}\left[\mathbf 1\{s_t>q_t\}-\alpha\right]
	=\sum_{t=1}^{T}\left(q_{t+1}-q_t\right)
	=q_{T+1}-q_1 .
	\]
	Dividing by $\eta T$ and applying Step 1,
	\[
	\left|\frac1T\sum_{t=1}^{T}\mathbf 1\{s_t>q_t\}-\alpha\right|
	=\frac{|q_{T+1}-q_1|}{\eta T}
	\le\frac{2R+\eta}{\eta T},
	\]
	which is equation~(16) of the manuscript. Substituting $\eta=R/\sqrt T$ gives
	$\eta T=R\sqrt T$ and the stated rate.
\end{proof}

\begin{remark}[What this argument replaces, and what it does not deliver]
	\label{rem:ogd_scope}
	A bound of this form is often obtained by invoking the
	online-convex-optimisation regret guarantee of \citet{zinkevich2003online}
	and converting average regret on the pinball loss into average miscoverage.
	That conversion requires a local Lipschitz argument in the indicator
	parametrisation and does not follow from the regret bound alone. The proof
	above needs none of it: the update telescopes exactly, and the only work is
	showing the iterates cannot escape a bounded interval. The constant is
	explicit and the statement is pathwise.

	Three limitations should be read alongside it. It controls the
	\emph{time-average} of the coverage indicator, not conditional coverage and
	not coverage over any subwindow: a sequence can satisfy
	equation~(16) of the manuscript while violating badly in one regime and
	over-covering in another, which is exactly the behaviour
	Section~6 of the manuscript warns about. It requires bounded scores, which
	rules out nothing empirically but is an assumption. And, most importantly
	for this paper, it is a statement about the \emph{online update}
	equation~(15) of the manuscript, whereas the rolling estimator used throughout the
	empirical sections is a trailing-window empirical quantile, which is not
	literal online gradient descent. In stationary regimes the two target the
	same population quantile; under drift the window trades bias for variance
	near transitions. We therefore cite this bound for what it is --- a
	dependence-free guarantee for one estimator in the family --- and not as
	cover for the estimator we run.
\end{remark}

\subsection{Proof of Proposition~4.8 (finite-window coverage under drift)}
\label{app:rolling_drift}

Proposition~4.8 of the manuscript is a supporting bound. Its conclusion holds
under the total-variation drift measure $\delta_w(t)$ of equation~(17) of the
manuscript, which is not estimable from a single path without further
assumptions, and the constant it produces is too loose to be informative at
$\alpha = 0.01$. It is retained because it makes precise which quantity has to be
small for a trailing window to work; $\hat\delta_w(t)$ below is the operational
monitoring diagnostic that replaces it in practice.


\begin{proof}[Proof of Proposition~4.8 of the manuscript]
	Let $S_t=\hat q_t^{lo}-r_t$ and let $\qVroll$ be the trailing-window
	conformal shift of equation~(9) of the manuscript, the
	$\lceil (w+1)(1-\alpha)\rceil$-th order statistic of
	$\{S_\tau\}_{\tau=t-w}^{t-1}$. It differs from the plain empirical quantile
	$Q_{1-\alpha}$ of the same window by one order statistic, an
	$\mathcal O(1/w)$ gap that is absorbed into $\Delta_w$ below without changing
	its rate.
	The coverage event satisfies
	\begin{equation}
		\label{eq:app_rolling_score_event}
		r_t\ge \hat q_t^{lo}-\qVroll
		\Longleftrightarrow
		S_t\le \qVroll.
	\end{equation}
	
	Following the weighted conformal framework of
	\citet{barber2023conformal}, the coverage deviation under
	non-exchangeability is bounded by the average total variation distance
	between calibration-window score distributions and the test distribution:
	\begin{equation}
		\label{eq:app_tv_bound}
		\left|
		\mathbb P(S_t\le\qVroll)-(1-\alpha)
		\right|
		\le
		\frac{1}{w}
		\sum_{\tau=t-w}^{t-1}
		d_{\mathrm{TV}}\!\left(\mathcal L(S_\tau),\mathcal L(S_t)\right).
	\end{equation}
	By definition,
	\[
	\frac{1}{w}
	\sum_{\tau=t-w}^{t-1}
	d_{\mathrm{TV}}\!\left(\mathcal L(S_\tau),\mathcal L(S_t)\right)
	\le
	\delta_w(t).
	\]
	The empirical quantile estimated from a dependent window of size $w$
	adds the finite-sample remainder
	\begin{equation}
		\label{eq:app_delta_w}
		\Delta_w\le C\sqrt{\frac{\log w}{w}}.
	\end{equation}
	Combining the drift and estimation terms yields
	\begin{equation}
		\label{eq:app_rolling_score_bound}
		\mathbb P(S_t\le\qVroll)
		\ge
		1-\alpha-\delta_w(t)-\Delta_w.
	\end{equation}
	Returning to the score representation gives
	\begin{equation}
		\label{eq:app_rolling_bound}
		\mathbb P\!\left(
		r_t\ge \hat q_t^{lo}-\qVroll
		\right)
		\ge
		1-\alpha-\delta_w(t)-\Delta_w.
	\end{equation}
\end{proof}

The bound separates distributional drift, captured by $\delta_w(t)$, from
finite-sample estimation error under dependence, captured by $\Delta_w$.
In the empirical implementation, $\delta_w(t)$ is proxied by the empirical
total variation distance between the first and second halves of the rolling
window. Using $\sqrt w$ histogram bins, define
\begin{equation}
	\label{eq:app_delta_hat}
	\hat\delta_w(t)
	=
	\frac{1}{2}\sum_b |\hat p_{1,b}-\hat p_{2,b}|,
\end{equation}
where $\hat p_{1,b}$ and $\hat p_{2,b}$ are empirical bin frequencies over
the first and second halves of the window. When $\hat\delta_w(t)$ exceeds a
practitioner-chosen threshold, the bound weakens and a shorter window or
temporary deployment pause may be warranted.

%============================================================

\subsection{Proof of Proposition~7.1 of the manuscript}
\label{app:identification}

Recall $u_t = F_t(\hat q^{lo}_t)$ and $V_t = \mathbf 1\{r_t < \hat q^{lo}_t\}$,
with $\hat q^{lo}_t$ measurable with respect to $\mathcal G_{t-1}$.

\begin{proof}
	Since $\hat q^{lo}_t$ is $\mathcal G_{t-1}$-measurable and $F_t$ is the
	conditional law of $r_t$ given $\mathcal G_{t-1}$,
	\[
	\P\!\left(V_t = 1 \mid \mathcal G_{t-1}\right)
	= \P\!\left(r_t < \hat q^{lo}_t \mid \mathcal G_{t-1}\right) = u_t ,
	\]
	and $u_t$ is $\mathcal G_{t-1}$-measurable. For $s \le t-1$ the indicator
	$V_s$ is a measurable function of $(r_s, \hat q^{lo}_s)$ and hence
	$\mathcal G_{t-1}$-measurable. Conditioning on $\mathcal G_{T-1}$,
	\[
	\P\!\left(V_{1:T}=v_{1:T}\right)
	= \E\!\left[\mathbf 1\{V_{1:T-1}=v_{1:T-1}\}\,
	u_T^{\,v_T}(1-u_T)^{1-v_T}\right],
	\]
	because the first factor is $\mathcal G_{T-1}$-measurable and the conditional
	law of $V_T$ is Bernoulli$(u_T)$. Iterating the same step down to $t=1$ gives
	the path likelihood stated as equation (12) of the manuscript.

	The right-hand side is a functional of the joint law of $(u_1,\dots,u_T)$
	alone. Hence if $P$ and $Q$ induce the same law for $\{u_t\}$ they induce the
	same law for $V_{1:T}$, so $\E_P[\phi(V)] = \E_Q[\phi(V)]$ for every
	measurable $\phi$, and in particular the rejection probability of any test in
	that class is the same under both: its power equals its size.

	For the moment-restriction case, $\E[V_t \mid \mathcal G_{t-1}] = u_t$, so
	$u_t \equiv \alpha$ implies $\E[(V_t-\alpha)Z_{t-1}] = 0$ for every
	predictable $Z_{t-1}$ with finite second moment. Two models both satisfying
	$u_t\equiv\alpha$ therefore satisfy every such restriction exactly, and a test
	of one of them rejects each with probability equal to its nominal size. The
	dynamic quantile regression of \citet{engle2004caviar} is of this form, with
	$Z_{t-1}$ collecting lagged exceedance indicators and the reported quantile
	itself.
\end{proof}

\begin{remark}[What the proposition does not cover]
	\label{rem:identification_scope}
	Two exclusions are worth stating. First, $\sigma(V_t, \hat q^{lo}_t)$ is
	strictly larger than $\sigma(V_t)$, and the two models constructed below do
	\emph{not} have the same marginal law for $\hat q^{lo}_t$. A diagnostic that
	compared the level of the reported quantile against a scale reference would
	separate them --- but that reference is not in $\sigma(V,\hat q^{lo})$, and no
	test in standard use within that class exploits the difference. This is why
	the proposition is stated for the moment-restriction family rather than for
	$\sigma(V,\hat q^{lo})$ in full. Second, statistics reading the
	\emph{magnitude} of an exceedance rather than its occurrence lie outside the
	class entirely; Supplement~\ref{app:es} treats them.
\end{remark}

\subsection{Proposition~\ref{prop:blindspot_depth}: the construction, and the depth of the blind spot}
\label{app:blindspot_depth}

Let $G$ be the standardised Student-$t_5$ law of the correctly specified model,
with $\alpha$-quantile $q_{\mathrm{true}} = \nGapQTrueStd$ in units of
conditional scale. A forecaster that truncates its predictive support, discarding
mass $\delta$ from each tail and renormalising, reports
\[
q_\delta \;=\; G^{-1}\!\left(\delta + \alpha(1-2\delta)\right),
\]
which is closer to zero than $q_{\mathrm{true}}$ and is what
$\texttt{top\_k}$ truncation does to a categorical predictive law.

For the two models to induce the same $u$-process it suffices that the
alternative return law $H$ place exactly $\alpha$ mass strictly below $q_\delta$,
since then $u_t \equiv \alpha$ under both. Requiring in addition that $H$ be
symmetric with unit variance and the same mean absolute deviation as $G$, and
that its density be non-increasing in $|x|$, gives a feasibility problem that is
\emph{linear} in the atoms of a discrete symmetric law on a fixed grid:
\begin{align*}
	\textstyle\sum_i p_i = 1, \qquad
	\textstyle\sum_i p_i x_i^2 = 1, \qquad
	\textstyle\sum_i p_i |x_i| = \E_G|X|, \\
	\textstyle\sum_{i:\,x_i < q_\delta} p_i = \alpha, \qquad
	p_{i+1} \le p_i \ \text{ for } x_i > 0, \qquad p_i \ge 0 .
\end{align*}
Feasibility is therefore decided by a linear programme rather than argued.
Bisection on $\delta$ gives the largest feasible truncation depth
$\delta^\star = \nGapDeltaUni$, at which the reported quantile is
$\nSupQClassUni$ and the Value-at-Risk is understated by \nGapUndUni\%.

\begin{table}[H]
	\centering
	\caption{$\delta^\star$ as a functional of the restriction placed on the
		alternative return law, at $\alpha = 0.01$ against a standardised
		Student-$t_5$ honest model. Grid spacing $\nSupDeltaGrid$, ceiling
		$\nSupDeltaCeiling$;
		$\delta^\star$ converges from above as the grid is refined, so each entry
		is an upper bound and the bias overstates the blind spot.}
	\label{tab:delta_by_class}
	\footnotesize
	\begin{tabular}{@{}lrrr@{}}
		\hline\hline
		Restriction on the return law & $\delta^\star$ & reported $q$ & understated \\
		\hline
		none & \nGapDeltaFree & $\nSupQClassFree$ & \nGapUndFree\% \\
		unimodal & \nGapDeltaUni & $\nSupQClassUni$ & \nGapUndUni\% \\
		unimodal, fourth moment $\le$ honest model's & \nSupDeltaClassMoment & $\nSupQClassMoment$ & \nSupUndClassMoment\% \\
		unimodal, Pareto tail index 5 beyond $3\sigma$ & \nSupDeltaClassPareto & $\nSupQClassPareto$ & \nSupUndClassPareto\% \\
		GARCH class, standardised Student-$t$ innovations & --- & $\nSupQClassGarchT$ & \nSupUndClassGarchT\% \\
		\hline\hline
	\end{tabular}
	\begin{minipage}{\linewidth}\scriptsize
		\smallskip
		The GARCH row carries no $\delta^\star$: that class fixes the innovation
		family to a one-parameter standardised Student-$t$, so truncation depth
		does not index it. The entry is the lightest $\alpha$-quantile the family
		can reach, the Gaussian limit.
	\end{minipage}
\end{table}

The resulting pair is matched to machine precision on coverage
($u \equiv 0.01$), variance, mean absolute deviation, median and predictive
standard deviation, and to three decimals on point-forecast mean absolute error
and root mean squared error. Both laws place exactly $\alpha$ mass strictly below
their own reported threshold, so the exceedance indicator is Bernoulli($\alpha$)
and serially independent under each, and every test of the exceedance path has
the same null distribution on the two. Simulated at $T = \nSupPairT$ the realised
rates are $\nSupPairPiHonest$ and $\nSupPairPiAlt$, both paths are Basel Green,
and neither rejects at any conventional level on Kupiec
($p = \nSupPairKupHonest$ and $\nSupPairKupAlt$), on Christoffersen independence
($\nSupPairCCHonest$ and $\nSupPairCCAlt$), or on a dynamic quantile test with
\nSupPairDQLags{} lagged exceedance indicators and the reported quantile as
instruments ($\nSupPairDQHonest$ and $\nSupPairDQAlt$). The simulation checks
that the sampled law reproduces the constraints the linear programme imposed;
the indistinguishability itself is the equality above. The reported thresholds
are $\nPairVaRHonest$ and $\nPairVaRAlt$ in units of conditional scale.

Two further properties are worth recording. Truncation is not repairable by the
sampling parameters an analyst would vary: restricting the support to its $k$
most probable atoms bounds every quantile read from any reweighting of that
support below by $\min S_k$, so no temperature and no nucleus threshold recovers
a quantile the truncation has removed. And matching \emph{mean} expected
shortfall as a fifth linear constraint is also feasible, but it does not make
magnitude-reading statistics blind: the law of $Z_2$ is not pinned by its mean,
and at $T = \nSupPowerT$ a size-$5\%$ one-sided test calibrated on the honest
model still rejects the alternative with probability $\nSupPowerRej$ over
\nSupPowerReps{} replications, a $95\%$ interval of $[\nSupPowerLo,
\nSupPowerHi]$ that excludes the nominal size. Necessary, not
sufficient.


\bibliographystyle{elsarticle-harv}
\bibliography{calibrating_the_oracle}
\end{document}
